%!TEX program = uplatex
\RequirePackage{plautopatch}
\documentclass[uplatex,dvipdfmx]{jlreq}
\usepackage{amsmath,amssymb}

\pagestyle{empty}
\begin{document}

\begin{center}
  {\Huge\bfseries 応用特異点論}
\end{center}

\vspace{10pt}

\begin{center}
  {\Large\bfseries 数学（厳密科学）としてのカタストロフ理論をめざして}
\end{center}

\vspace{10pt}

\begin{center}
  {\large 前文}
\end{center}

\medskip

泉屋・石川著「応用特異点論」（共立、1998年）序文より抜粋：

「万物は特異点である」と言った古代の哲学者がいたかどうかはわからないが、
人類は古代より「特異的」なものに注目してきた。
数学においても、考えてみれば、正三角形、球、平面などちょっと歪めれば、
他のものに変化してしまうような危うい対象を美しいとみて研究の主要な対象としてきた。
さらに、われわれが物を認識する場合、周りの状況と異なっている部分を認識している。
たとえば、物を見る場合、その輪郭やひだを見ることにより形を理解している。
これらは、周囲と異なるものということで、特異点の集まりであるといえる。

「応用特異点論」は広く認められている分野の名前ではない。
いうなれば我々の造語である。……中略……
応用特異点論を応用する対象は数学以外の諸科学のみならず、
数学の内部のさまざまな分野に及ぶ。
その意味では、この分野は応用・純粋双方の数学にまたがる分野であるといえる。

\end{document}